\begin{lema}[Reduction to rooks on the last board]\label{terminal-pivotes}
Let $A$ be the mask of the board $B^{(K)}$. On $F\cap D_K^c$ there is a placement completing the
board, and
\[
 Q(B^{(K)})=\per A,\qquad L(B^{(K)})=L_T(A),
\]
where $L_T(A)=H_T(p)+\log\Gamma_2$ is the rook predictor on the same support, as in the
masked rook model of~\cite[Sections~8.1 and~8.5]{pred}. Here
$H_T(p)=-\sum_e p_e\log p_e+\sum_e(1-p_e)\log(1-p_e)$ is the rook entropy, evaluated at its
stationary point with row and column masses equal to one, and $\Gamma_2$ is the factor from
Section~\ref{costes:permanente}.
\end{lema}
\begin{proof}
On $F$ the Bethe point $p$ of $B^{(K)}$ has row and column masses equal to one and is positive on all
available squares. For any set $I$ of rows, its entire mass $|I|$ goes to the set of neighboring
columns, which admits at most one unit per column. Consequently $|\mathcal N(I)|\ge|I|$.
To see directly that every row can be covered, take a matching of maximum cardinality. If it
leaves a row uncovered, start from all uncovered rows and follow alternating paths, taking
unmatched edges from rows and matched edges from columns. A reachable uncovered column would give
an augmenting path, a contradiction. Every reachable column is therefore matched to a reachable
row, and these matched rows do not exhaust the reachable rows, which include an uncovered row.
Every neighbor of a reachable row is reachable, including the column from which that row was
entered, if any. This contradicts the inequality for the set of reachable rows. Thus $\per A>0$.

On $D_K^c$ each diagonal contains at most one empty crossing. Every permutation on the support
therefore satisfies the constraints from both diagonal families. Moreover, squares on occupied
diagonals were already excluded when forming $A$, so every permutation on the support gives a
legal completion of $B^{(K)}$. The construction is reversible and proves $Q(B^{(K)})=\per A$.

On a free diagonal with one available square $e$, the vacancy is $h_f=1-p_e=:q_e$. Each square
lies on two such diagonals, and a free diagonal with no squares has $h_f=1$. Thus in $H_B$ the
diagonal terms subtract two copies of $\sum_e q_e\log q_e$, leaving exactly the rook entropy
\[
 H_T(p)=-\sum_e p_e\log p_e+\sum_e q_e\log q_e .
\]
The stationarity equation of Lemma~\ref{coste-vacancias} becomes $p_eq_e=a_ib_j$ on every
available square, with the same $a_i,b_j$. The equality $H_B=H_T$ also follows directly as an
identity of functionals at the feasible points on the support, since on every single-square
diagonal the vacancy is always the complement of the sole entry. Hence the Bethe point is also
the stationary point of the rook predictor.

It remains to verify the normalizer, as in \cite[Section~8.1]{pred}. For a square of mass $p$,
put $q=1-p$, $\tau_d=pq/(1-dp)$, and $b_d=(1-dp)/q^d$. When we eliminate the coordinate of its
first diagonal, the pivot in $S_4$ is $q+\tau_4$. If $r$ is the indicator vector of the other lines
of the square,
the Schur complement replaces the contribution $\tau_4rr^T$ in the remaining block by
$(\tau_4-\tau_4^2/(q+\tau_4))rr^T$. We then eliminate the second diagonal of that square with
$d=3$. At both steps,
\[
 q+\tau_d=\frac{q(1-(d-1)p)}{1-dp},\qquad
 \tau_d-\frac{\tau_d^2}{q+\tau_d}=\tau_{d-1},\qquad
 b_d(q+\tau_d)=b_{d-1}.
\]
Since the eliminated diagonal contains a single square, the operation affects only that square's
contribution. A line with no squares has pivot one and does not affect the other coordinates.
The zero mode $(\mathbf1_R,-\mathbf1_C,0)$ has zero coordinate on every eliminated line. We may
therefore calculate the pseudodeterminant in an orthonormal basis of its orthogonal complement
that includes the coordinate vectors of the eliminated diagonals. In this basis the ordinary
factorization using the pivots is exact, and the row and column mode retains its length
$\sqrt{2K}$. The corresponding factor in the normalizer is $\sqrt{2\pi}\sqrt{2K}=\sqrt{4\pi K}$
both before and after elimination. Since $e^{C_d}=\prod_e b_d(p_e)$, the identities prove
$e^{C_4}\pdet S_4=e^{C_2}\pdet S_2$ after the two eliminations per square. In particular,
$\Gamma_4=\Gamma_2$. This gives $L(B^{(K)})=H_T(p)+\log\Gamma_2=L_T(A)$.
\end{proof}
